% Options for packages loaded elsewhere
\PassOptionsToPackage{unicode}{hyperref}
\PassOptionsToPackage{hyphens}{url}
\documentclass[
]{article}
\usepackage{xcolor}
\usepackage{amsmath,amssymb}
\setcounter{secnumdepth}{-\maxdimen} % remove section numbering
\usepackage{iftex}
\ifPDFTeX
  \usepackage[T1]{fontenc}
  \usepackage[utf8]{inputenc}
  \usepackage{textcomp} % provide euro and other symbols
\else % if luatex or xetex
  \usepackage{unicode-math} % this also loads fontspec
  \defaultfontfeatures{Scale=MatchLowercase}
  \defaultfontfeatures[\rmfamily]{Ligatures=TeX,Scale=1}
\fi
\usepackage{lmodern}
\ifPDFTeX\else
  % xetex/luatex font selection
\fi
% Use upquote if available, for straight quotes in verbatim environments
\IfFileExists{upquote.sty}{\usepackage{upquote}}{}
\IfFileExists{microtype.sty}{% use microtype if available
  \usepackage[]{microtype}
  \UseMicrotypeSet[protrusion]{basicmath} % disable protrusion for tt fonts
}{}
\makeatletter
\@ifundefined{KOMAClassName}{% if non-KOMA class
  \IfFileExists{parskip.sty}{%
    \usepackage{parskip}
  }{% else
    \setlength{\parindent}{0pt}
    \setlength{\parskip}{6pt plus 2pt minus 1pt}}
}{% if KOMA class
  \KOMAoptions{parskip=half}}
\makeatother
\setlength{\emergencystretch}{3em} % prevent overfull lines
\providecommand{\tightlist}{%
  \setlength{\itemsep}{0pt}\setlength{\parskip}{0pt}}
\usepackage{bookmark}
\IfFileExists{xurl.sty}{\usepackage{xurl}}{} % add URL line breaks if available
\urlstyle{same}
\hypersetup{
  pdftitle={Notes},
  hidelinks,
  pdfcreator={LaTeX via pandoc}}

\title{Notes}
\author{}
\date{}

\begin{document}
\maketitle

86\\
LOOP := Menge aller LOOP Funktionen

103:

\begin{itemize}
\tightlist
\item
  Falls {}
\item
  fuer {} gilt\\
  - innere While\\
  - wirkt wie abrundende division mit 2\\
  - if\\
  - testet ungerade\\
  - wurde zugewiesen rundet?\\
  - = war {} ungerade =\textgreater{} x = 6*x + 4 zugewiesen\\
  - aeussere Schleife\\
  - Falls x gerade =\textgreater{} x/2\\
  - sonst =\textgreater{} x = 3*x+1 (*)\\
  f(a):\\
  - 0 falls a \textgreater= 1\\
  - 0 falls a \textgreater{} 1 und (x) zu x = 1 fuehrt\\
  - sonst n.d.\\
  Gilt f(a) = 0 fuer alle {}?\\
  Vermutung: ja - wurde aber noch nicht bewiesen/widerlegt
\end{itemize}

\begin{quote}
Collatz Problem v. 1937
\end{quote}

Fazit: Es kann schwierig sein, die von einem Programm berechnete
Funktion zu bestimmen.\\
Insbesondere die Equivalence von Programmen ist schwer zu testen
=\textgreater{} ist f(a) zu g(x) = 0 äquivalent

\subsection{Registermaschinen}\label{registermaschinen}

119

\end{document}
